% % Guard: macros.tex is \input both by main.tex (after the packages it needs,
% e.g. amsmath for \DeclareMathOperator) and by individual section files so
% they remain compilable on their own. Loading it twice would redefine every
% \newcommand and abort, so the second and later loads are no-ops.
\ifdefined\PQRmacrosloaded\endinput\fi
\let\PQRmacrosloaded\relax

% --- 0. revision marking ---
% Text added to Section 2 during the current revision is wrapped in \added{...}
% so the advisors can see at a glance what is new. Redefine to \relax to print
% the paper without the marking.
\newcommand{\added}[1]{\textcolor{blue}{#1}}
\long\def\added#1{\begingroup\color{blue}#1\endgroup}  % \long: may span paragraphs


% --- A. operator-style score and function names ---
\newcommand{\HD}{\mathit{HD}}                        % histogram distance; use \HD_\gamma
\newcommand{\Attr}{\mathit{Attr}}                    % attribute set of a schema
\newcommand{\Attrs}{\mathit{Attrs}}                  % attribute set argument
% Abbreviated in MATH only; the prose still spells the names out in full.
% Spelling them out inside formulas makes displayed equations far too wide
% for a two-column layout (the product-rule step below overflowed by 200pt).
\newcommand{\Use}{\mathit{Use}}
\newcommand{\Cov}{\mathit{Cov}}
\newcommand{\Rep}{\mathit{Rep}}
\newcommand{\DivS}{\mathit{Div}}
\newcommand{\TVD}{\mathit{TVD}}
\newcommand{\GS}{\mathit{GS}}                        % global sensitivity
\DeclareMathOperator{\Dom}{Dom}
\DeclareMathOperator{\Lap}{Laplace}

% --- B. core objects ---
\newcommand{\Qone}{Q_1}
\newcommand{\Qtwo}{Q_2}
\newcommand{\pD}{p(D)}
\newcommand{\views}{\mathcal{P}}
\newcommand{\hist}[2]{H_{#1,#2}}                     % \hist{\Qone(\pD)}{A}
\newcommand{\dhist}{\Delta H}
\newcommand{\nhist}{\widetilde{\Delta H}}            % the released noisy object

% --- C. privacy budget and parameters ---
\newcommand{\eps}{\epsilon}
\newcommand{\epsP}{\epsilon_P}
\newcommand{\epsS}{\epsilon_S}
\newcommand{\epsH}{\epsilon_H}
\newcommand{\gammaref}{\gamma_{\mathrm{ref}}}
\newcommand{\gammaeval}{\gamma_{\mathrm{eval}}}
\newcommand{\taumax}{\tau_{\max}}

% --- D. mechanisms ---
\newcommand{\ExpMech}{\textsc{ExpMech}}
\newcommand{\OneShotTop}[1]{\textsc{OneShotTop-}#1}
\newcommand{\AdaptiveTau}{\textsc{AdaptiveTau}}
% --- E. experiment methods (used in sections/experiments.tex); a method is renamed here ---
% (the figure legends are drawn by src/pqde/paper_figures.py and must be redrawn after a rename)
\newcommand{\PrivExpl}{\textsf{PrivExpl}}
\newcommand{\OptPrivExpl}{\textsf{OptPrivExpl}}
\newcommand{\BF}{\textsf{BF}}
\newcommand{\SingleStage}{\textsf{SingleStage}}
\newcommand{\RandomExpl}{\textsf{Random}}
\newcommand{\Bone}{\ensuremath{\mathsf{B1}}}
\newcommand{\Btwo}{\ensuremath{\mathsf{B2}}}
\newcommand{\Bthree}{\ensuremath{\mathsf{B3}}}

\section{Preliminaries}
\label{sec:preliminaries}

We present the data model and the notions of differential privacy that the paper relies on.

\subsection{Data, Queries and Views}
\label{subsec:prelim-db-queries}

\paragraph{Databases and histograms.}
Let $D$ be a relation with attribute set $\Attr(D)$~\cite{Foundations}.
Each tuple $t\in D$ is the record of one individual, and we write $N=|D|$.
Every attribute $A$ has a domain $\Dom(A)$, and $t[A]\in\Dom(A)$ denotes the value of $t$ on $A$.
For ease of presentation, we treat all attributes as categorical; a numerical attribute is bucketized into ranges, which then form its domain.
The domains are declared by the schema and are therefore \emph{public}: they do not depend on the contents of $D$, whereas a domain read off the data would itself disclose information about it. The \emph{histogram} of an attribute $A$ over a set of tuples $T\subseteq D$ is the function $H_{T,A}:\Dom(A)\rightarrow\mathbb{N}$ with $H_{T,A}(v)=|\{t\in T \mid t[A]=v\}|$. For $T \neq \emptyset$, the \emph{distribution} of $A$ over $T$ is
$P_{T,A}(v) = H_{T,A}(v)/|T|$.

% \begin{definition}[Histogram]
% \label{def:histogram}
% For a set of tuples $T\subseteq D$ and an attribute $A \in \Attr(D)$, the \emph{histogram} of $A$ over $T$ is the function $H_{T,A}: \Dom(A)\rightarrow \mathbb{N}$ with
% \[
% H_{T,A}(v) = |\{t\in T \mid t[A]=v\}| .
% \]
% \end{definition}

\paragraph{Queries.}
We consider queries that select tuples of $D$ by a Boolean predicate and possibly aggregate them,
\[
\texttt{SELECT } agg(A)
\texttt{ FROM } D
\texttt{ WHERE } \varphi ,
\]
where $agg \in \{\texttt{COUNT},\texttt{SUM},\texttt{AVG}\}$ and $\varphi$ is a Boolean combination of \emph{atomic conditions}: comparisons such as $A \leq c$ or $A \geq c$, equalities $A = v$, and set memberships $A \in S$.
Our explanations describe the \emph{population} that a query selects rather than the scalar that it may return, since it is the population whose composition changes when a query is refined.
We therefore identify a query $Q$ with its selection and write
\[
Q(D) \;=\; \sigma_{\varphi}(D) \;=\; \{t \in D \mid \varphi(t)=\texttt{true}\}
\]
for the \emph{output} of $Q$ on $D$.

\paragraph{Views.}
Explanation frameworks commonly localize a change to a subpopulation that is defined by a predicate~\cite{FEDEX,TaoGMR22,wu2013}, and so do we.
A \emph{view predicate} $p$ is a conjunction of at most $\ell$ atomic conditions over a designated set of view attributes, for a small $\ell$ (at most $2$ in our experiments), and its \emph{view} is $p(D)=\{t\in D \mid p(t)=\texttt{true}\}$.
The framework is given a finite collection $\views$ of view predicates, enumerated from the public schema, so $\views$ does not depend on the contents of $D$.
It is also given two public sets of attributes, which may overlap: the \emph{explanation attributes} $\Attrs\subseteq\Attr(D)$, whose histograms may be released, and the \emph{reference attributes} $\mathcal{A}_R\subseteq\Attr(D)$, on which a view is compared with the whole database (Section~\ref{sec:explanation-quality}).
% The two sets are chosen independently and may overlap or coincide, as they do in our experiments; 
An attribute that a view predicate constrains is used neither as an explanation attribute nor as a reference attribute for that view.

Queries and views are both selections over $D$, so evaluating a query inside a view amounts to intersecting the two:
\begin{equation}
Q(p(D)) \;=\; \sigma_{\varphi\wedge p}(D) \;=\; Q(D)\cap p(D)\;\subseteq\; p(D).
\label{eq:view-restriction}
\end{equation}
In Example~\ref{ex:intro}, for instance, the view of $p = (\texttt{education}=\texttt{Masters})$ holds the Master's degree holders, and $\Qone(\pD)$ those among them who work at least $50$ hours a week.

\subsection{Differential Privacy}
\label{subsec:prelim-dp}

Differential privacy~\cite{DMNS2006,TheAlgFoundDP2014} requires a randomized mechanism to behave almost identically on any two databases that differ in one tuple. Let $\mathcal{D}$ denote the set of all databases over the schema of $D$.

\begin{definition}[Neighboring databases~\cite{DMNS2006}]
\label{def:neighboring}
Two databases $D$ and $D'$ are \emph{neighboring}, written $D\sim D'$, if $D'=D\cup\{t\}$ or $D'=D\setminus\{t\}$ for some tuple~$t$.
\end{definition}

We adopt this \emph{unbounded} form of the relation throughout, under which neighboring databases differ in size.

\begin{definition}[Differential privacy~\cite{DMNS2006,TheAlgFoundDP2014}]
\label{def:dp}
A randomized mechanism $\mathcal{M}$ satisfies $\epsilon$-differential privacy if for every pair of neighboring databases $D\sim D'$ and every set $S$ of outputs,
\[
\Pr[\mathcal{M}(D)\in S]
\leq
\exp(\epsilon)\,
\Pr[\mathcal{M}(D')\in S].
\]
\end{definition}

A smaller $\epsilon$ gives stronger privacy at the cost of more noise, which is calibrated to the sensitivity of what is computed.

\begin{definition}[Global sensitivity~\cite{TheAlgFoundDP2014}]
\label{def:global-sensitivity}
The \emph{global sensitivity} of a function $f:\mathcal{D}\rightarrow\mathbb{R}^d$ is
\[
\Delta_f
=
\max_{D\sim D'}
\|f(D)-f(D')\|_1 .
\]
We write $\Delta_f(N)$ for the same maximum over the neighboring pairs with $|D|=N$, so that $\Delta_f=\sup_N\Delta_f(N)$.
\end{definition}

\begin{definition}[Laplace mechanism~\cite{DMNS2006,TheAlgFoundDP2014}]
\label{def:laplace-mechanism}
Given a function $f:\mathcal{D}\rightarrow\mathbb{R}^d$ with global sensitivity $\Delta_f$, the Laplace mechanism releases $f(D)+\left(Y_1,\ldots,Y_d\right)$, where $Y_1,\ldots,Y_d$ are drawn independently from the zero-mean Laplace distribution $\Lap(\Delta_f/\epsilon)$ of scale $\Delta_f/\epsilon$.
\end{definition}

\begin{definition}[Exponential mechanism~\cite{mcsherry2007,TheAlgFoundDP2014}]
\label{def:exponential-mechanism}
Given a set $Y$ of candidates and a utility function $u:Y\times \mathcal{D}\rightarrow\mathbb{R}$ with global sensitivity $\Delta_u$, the exponential mechanism selects a candidate $y\in Y$ with probability proportional to $\exp\big(\epsilon\, u(y,D)/(2\Delta_u)\big)$.
\end{definition}

Both mechanisms satisfy $\epsilon$-differential privacy, and both are more accurate for functions of lower sensitivity.
They compose as follows.

\begin{theorem}[Composition and post-processing~\cite{TheAlgFoundDP2014,DBLP:series/synthesis/2016Li}]
\label{thm:sequential-composition}\label{thm:post-processing}
Let $\mathcal{M}_1,\ldots,\mathcal{M}_k$ be mechanisms with $\mathcal{M}_i$ satisfying $\epsilon_i$-DP.
\begin{enumerate}[label=(\roman*),leftmargin=*]
\item \emph{Composition.} Releasing $(\mathcal{M}_1(D),\ldots,\mathcal{M}_k(D))$ satisfies $(\sum_{i}\epsilon_i)$-DP; the same bound holds in the \emph{adaptive} setting, where the choice of $\mathcal{M}_i$ may depend on the outputs of $\mathcal{M}_1,\dots,\mathcal{M}_{i-1}$.
\item \emph{Post-processing.} For any randomized $f$ that reads $D$ only through $\mathcal{M}_1(D)$, the mechanism $f(\mathcal{M}_1(D))$ is also $\epsilon_1$-DP.
\end{enumerate}
\end{theorem}


We use the adaptive form, since the utility of each round of Algorithm~\ref{alg:topk} depends on the explanations released in earlier rounds.
